\begin{figure} [h]
	\centering
	\pgfplotstableread[col sep=semicolon] {misc/chap_4_ball_bouncing/perturbations_significance_post_hoc.csv}\datatable
	\tikzsetnextfilename{perturbations_significance_post_hoc}
	\begin{tikzpicture}
		
		\begin{axis}[%
			height=0.3\linewidth,
			scale only axis,
			%xmin=0.65,
			%xmax=2.35,
			%xtick={1,1.5,2},
			xtick={0,1,2,3},
			xticklabels={{$c_{pre}$},{$c_{post}^1$},{$c_{post}^2$},{$c_{post}^3$}},
			%ymin=1.06276752948761,
			%ymax=1.78985784053802,
			ylabel=$\mean{\epsilon_r}$ (\si{\centi\meter}),%$\sigma(\epsilon_r)$
			ylabel near ticks,
			yticklabel={\pgfmathparse{\tick*100}\pgfmathprintnumber{\pgfmathresult}}, % conversion from m to cm
			scaled ticks=false, 
			tick label style={/pgf/number format/fixed},
			ylabel style={yshift=-1mm},
			%boxplot/draw direction=y, 
			]
			\addplot [color=Orient] table[x=source, y=mean_avg]{\datatable};
			\addplot [color=Orient, only marks, mark=*, mark size=1.7pt, mark options={fill=white, fill opacity=1}]
			plot [error bars/.cd, y dir = both, y explicit]
			table[x=source, y=mean_avg, y error=SEM_avg]{\datatable};		
		\end{axis}		
	\end{tikzpicture}%
	\caption{Effets des perturbations sur $\mean{\epsilon_r}$, les erreurs types sont utilisées comme intervalle de confiance.}
	\label{fig_perturbations_significance_post_hoc}
\end{figure} 